\documentclass[11pt,a4paper]{article}

\usepackage[margin=1in]{geometry}
\usepackage{amsmath,amssymb,amsthm}
\usepackage{array}
\usepackage[colorlinks=true,linkcolor=blue,citecolor=blue,urlcolor=blue]{hyperref}

\theoremstyle{plain}
\newtheorem{theorem}{Theorem}[section]
\newtheorem{proposition}[theorem]{Proposition}
\newtheorem{lemma}[theorem]{Lemma}
\newtheorem{corollary}[theorem]{Corollary}

\theoremstyle{definition}
\newtheorem{definition}[theorem]{Definition}
\newtheorem{example}[theorem]{Example}

\theoremstyle{remark}
\newtheorem{remark}[theorem]{Remark}

\newcommand{\Cl}{\mathrm{Cl}}
\newcommand{\CCl}{\mathbb{C}\mathrm{l}}
\newcommand{\End}{\mathrm{End}}
\newcommand{\Hom}{\mathrm{Hom}}
\newcommand{\rad}{\mathrm{rad}}
\newcommand{\Span}{\mathrm{span}}
\newcommand{\Tr}{\mathrm{Tr}}
\newcommand{\id}{\mathrm{id}}

\title{\textbf{Spinors as Minimal Left Ideals: Exact Split Characterization and Structural Refutations of Conjecture 00000007604}}
\author{Formal Resolution and Verification Report}
\date{November 2026}

\begin{document}
\maketitle

\begin{abstract}
We resolve Conjecture \texttt{00000007604}, which defines algebraic spinors as elements of minimal left ideals of a Clifford algebra and asserts that (i) the realization dimension equals the half-dimension power of two $2^{n/2}$, (ii) the spectrum is complete, (iii) the carrier of the realization is the minimal ideal, and (iv) the carrier of the minimal ideal is the matrix realization of the simple module.
We prove that as an unqualified universal statement for quadratic spaces $(V, Q)$ of dimension $n = \dim_K V$, the conjecture is \textbf{false}, failing across three independent structural regimes:
\begin{enumerate}
  \item \textbf{Odd dimensions ($n = 2m + 1$):} $2^{n/2} = 2^m\sqrt{2}$ is irrational, so no module over any field can have dimension $2^{n/2}$; moreover, in split/complex odd Clifford algebras $\Cl_{m+1,m}(K) \cong M_{2^m}(K) \oplus M_{2^m}(K)$ (such as $\Cl_{1,0}(K) \cong K \times K$ at $n = 1$), every minimal nonzero left ideal has dimension $2^{(n-1)/2} \neq 2^{n/2}$, fails to be faithful, and realizes only half of the central/generator eigenvalue spectrum.
  \item \textbf{Even-dimensional division-algebra obstructions (real/rational Bott periodicity):} Even when $n = 2m$ is even and $Q$ is nondegenerate, if the Wedderburn division algebra $\mathbb{D}$ is nontrivial, minimal left ideals have dimension $(\dim_K \mathbb{D})^{1/2} \cdot 2^{n/2} > 2^{n/2}$. Concretely, in $\Cl_{0,2}(K) \cong \mathbb{H}_K$ ($n = 2$, $K \in \{\mathbb{Q}, \mathbb{R}\}$), every nonzero element is invertible, the unique minimal nonzero left ideal is the full $4$-dimensional algebra ($\dim_K S = 4 \neq 2 = 2^{2/2}$), left multiplication embeds $\mathbb{H}_K$ properly into $M_4(K)$ rather than $M_2(K)$, and Clifford generators have no $K$-eigenvalues.
  \item \textbf{Degenerate quadratic forms ($Q$ degenerate, e.g.\ $Q = 0$):} Over \emph{every} field $K$ (including $\mathbb{C}$) and in even dimension $n = 2m \ge 2$, the exterior Clifford algebra $\Cl(V, 0) = \Lambda(V)$ is a local non-semisimple ring whose unique minimal nonzero left ideal is the $1$-dimensional top exterior line $\Lambda^n(V)$ ($\dim_K S = 1 \neq 2^{n/2}$), which satisfies $S^2 = 0$, contains no nonzero idempotent, and carries the trivial zero action of $V$.
\end{enumerate}
Conversely, we prove a positive non-vacuity and exact characterization theorem: under the classical Cartan--Chevalley reading of ``minimal ideal'' as ``minimal left ideal,'' all four clauses hold simultaneously whenever $(V, Q)$ is an \emph{even-dimensional nondegenerate split (neutral) quadratic space} ($\Cl(V, Q) \cong M_{2^m}(K)$, which over $\mathbb{C}$ covers all nondegenerate forms of even dimension $n = 2m$, and over $\mathbb{R}$ covers signatures $p + q = 2m$ with $p - q \equiv 0, 2 \pmod{8}$ under $v^2 = +Q(v)\cdot 1$). Both the positive witness $\Cl_{1,1}(\mathbb{Q}) \cong M_2(\mathbb{Q})$ and complete minimal-ideal classifications for the counterexamples $\Cl_{1,0}(\mathbb{Q})$, $\Cl_{0,2}(\mathbb{Q}) \cong \mathbb{H}_{\mathbb{Q}}$, and $\Lambda(\mathbb{Q}^2)$ are machine-checked in Core Lean~4.19.0 with zero external dependencies and zero \texttt{sorry}s.
\end{abstract}

\section{Introduction and Analysis of the Problem Statement}

\subsection{Original Problem Statement}
Conjecture \texttt{00000007604} in \texttt{SOURCE.md} is stated bilingually in English and Chinese. The English text reads verbatim:
\begin{quote}
\textbf{Definition:} Spinors as elements of minimal left ideals.
\textbf{Conjecture:} The realization dimension equals the half-dimension power of two and the spectrum is complete; the carrier of the realization is the minimal ideal; and the carrier of the minimal ideal is the matrix realization of the simple module. \emph{(spinor half-dimension power of two)}
\end{quote}
The accompanying Chinese text carries the exact clause-by-clause correspondence:
\begin{enumerate}
  \item \emph{Definition:} Algebraic spinors are defined as elements of minimal (nonzero) left ideals $S \subset \Cl(V, Q)$ of the Clifford algebra associated with an $n$-dimensional quadratic vector space $(V, Q)$ over a base field $K$ ($\mathrm{char}\,K \neq 2$).
  \item \emph{Clause C1 (Half-dimension power of two):} The dimension of the spinor realization satisfies
  \begin{equation}\label{eq:half-dim}
    \dim_K S = 2^{n/2}, \qquad \text{equivalently} \qquad (\dim_K S)^2 = 2^n = \dim_K \Cl(V, Q).
  \end{equation}
  \item \emph{Clause C2 (Complete spectrum):} The spinor realization exhibits spectral completeness---both in the Cartan weight-space sense (joint eigenvalues of a maximal commuting family of Clifford involutions realize all $2^{n/2}$ sign patterns $(\pm 1, \dots, \pm 1)$ with multiplicity~$1$) and in the Burnside/Wedderburn operator sense (left multiplication $\rho_S : \Cl(V, Q) \to \End_K(S)$ is faithful and surjective onto the full endomorphism algebra).
  \item \emph{Clause C3 (Minimal ideal carrier):} The underlying vector space (carrier) of the spinor realization is the minimal [left] ideal $S$ itself.
  \item \emph{Clause C4 (Matrix realization of the simple module):} Under an algebra isomorphism $\Cl(V, Q) \cong M_{2^{n/2}}(K)$ sending a primitive idempotent $e \in \Cl(V, Q)$ to a diagonal matrix unit $E_{jj}$, the minimal left ideal $S = \Cl(V, Q) e$ is realized concretely as the $j$-th single-column subspace of $M_{2^{n/2}}(K)$, giving the standard column-vector matrix realization of the unique simple $\Cl(V, Q)$-module.
\end{enumerate}

\begin{remark}[Minimal left ideal vs.\ minimal two-sided ideal in Clauses C3--C4]\label{rem:two-sided}
Notice that the Definition explicitly specifies \emph{minimal left ideals}, whereas Clauses C3--C4 abbreviate this to ``the minimal ideal.'' Under the charitable reading (standard in Cartan~\cite{Cartan1966}, Chevalley~\cite{Chevalley1954}, and Lounesto~\cite{Lounesto2001}), ``the minimal ideal'' refers anaphorically to the minimal \emph{left} ideal introduced in the Definition; under this reading, Clauses C1--C4 form a celebrated true theorem in the even-dimensional split/complex nondegenerate regime (Theorem~\ref{thm:positive-split}). If instead ``minimal ideal'' were read literally as a minimal nonzero \emph{two-sided} ideal $J \trianglelefteq \Cl(V, Q)$, then even for $\Cl_{1,1}(K) \cong M_2(K)$ ($n = 2$), simplicity of $M_2(K)$ forces $J = M_2(K)$ of dimension $4 = 2^n \neq 2 = 2^{n/2}$, which as a left module decomposes as $S \oplus S$ and is \emph{not} simple. Even granting the charitable minimal-left-ideal reading throughout, Conjecture \texttt{00000007604} fails outside the even-dimensional split/complex nondegenerate class.
\end{remark}

\subsection{Summary of Results across Standard Clifford Settings}
Table~\ref{tab:overview} summarizes the status of Conjecture \texttt{00000007604} across the principal families of Clifford algebras $\Cl(V, Q)$ (using the convention $v^2 = +Q(v)\cdot 1$).

\begin{table}[ht]
\centering
\footnotesize
\setlength{\tabcolsep}{4.5pt}
\renewcommand{\arraystretch}{1.22}
\begin{tabular}{@{}llllll@{}}
\hline
\textbf{Regime / Example} & $\dim V$ & \textbf{Algebra $\Cl(V,Q)$} & $\dim_K S$ & $2^{n/2}$ & \textbf{Status} \\
\hline
Split even $\Cl_{m,m}(K)$, $\CCl_{2m}$ & $2m$ & $M_{2^m}(K)$ & $2^m$ & $2^m$ & \textbf{Holds} (Thm.~\ref{thm:positive-split}) \\
Real $\Cl_{p,q}(\mathbb{R})$, $p{-}q \equiv 0,2\;(8)$ & $2m$ & $M_{2^m}(\mathbb{R})$ & $2^m$ & $2^m$ & \textbf{Holds} (Thm.~\ref{thm:bott}) \\
\hline
Split/complex odd $\Cl_{m+1,m}(K)$, $\CCl_{2m+1}$ & $2m{+}1$ & $M_{2^m}(K) \oplus M_{2^m}(K)$ & $2^m$ & $2^m\sqrt{2}$ & \textbf{Fails} (Thm.~\ref{thm:odd-split}) \\
Concrete odd $\Cl_{1,0}(\mathbb{Q})$ & $1$ & $\mathbb{Q} \times \mathbb{Q}$ & $1$ & $\sqrt{2}$ & \textbf{Fails} (Cor.~\ref{cor:cl10}) \\
Non-split odd $\Cl_{0,1}(\mathbb{R})$ & $1$ & $\mathbb{C}$ & $2$ & $\sqrt{2}$ & \textbf{Fails} (Thm.~\ref{thm:bott}) \\
Quaternionic even $\Cl_{0,2}(\mathbb{Q}), \Cl_{0,2}(\mathbb{R})$ & $2$ & $\mathbb{H}_K$ (division) & $4$ & $2$ & \textbf{Fails} (Thm.~\ref{thm:cl02}) \\
Real $\Cl_{p,q}(\mathbb{R})$, $p{-}q \equiv 4,6\;(8)$ & $2m$ & $M_{2^{m-1}}(\mathbb{H})$ & $2^{m+1}$ & $2^m$ & \textbf{Fails} (Thm.~\ref{thm:bott}) \\
Exterior / degenerate $\Lambda(K^n)$ ($n \ge 1$) & $n$ & $\Lambda(K^n)$ (local) & $1$ & $2^{n/2}$ & \textbf{Fails} (Thm.~\ref{thm:exterior}) \\
\hline
\end{tabular}
\caption{Exact minimal left ideal dimension $\dim_K S$ versus $2^{n/2}$ across Clifford algebra regimes under $v^2 = +Q(v)\cdot 1$.}
\label{tab:overview}
\end{table}

\section{Algebraic Foundations}

Let $K$ be a field with $\mathrm{char}\,K \neq 2$, let $V$ be an $n$-dimensional $K$-vector space ($n \ge 0$), and let $Q : V \to K$ be a quadratic form with associated symmetric bilinear form
\[
  B_Q(u, v) = \frac{1}{2}\bigl(Q(u + v) - Q(u) - Q(v)\bigr), \qquad B_Q(v, v) = Q(v).
\]

\begin{definition}[Clifford algebra]
The \emph{Clifford algebra} $\Cl(V, Q)$ is the quotient of the tensor algebra $T(V) = \bigoplus_{k \ge 0} V^{\otimes k}$ by the two-sided ideal generated by $\{v \otimes v - Q(v)\cdot 1 : v \in V\}$. Equivalently, for every $u, v \in V \subset \Cl(V, Q)$,
\begin{equation}\label{eq:cliff-rel}
  v^2 = Q(v)\cdot 1, \qquad uv + vu = 2 B_Q(u, v)\cdot 1.
\end{equation}
If $(e_1, \dots, e_n)$ is a $B_Q$-orthogonal basis of $V$ with $Q(e_i) = q_i \in K$, then $e_i^2 = q_i \cdot 1$, $e_i e_j = -e_j e_i$ for $i \neq j$, and the ordered monomials $\{e_{i_1}\cdots e_{i_k} : 1 \le i_1 < \dots < i_k \le n\}$ form a $K$-basis of $\Cl(V, Q)$ of dimension $\dim_K \Cl(V, Q) = 2^n$ (Poincar\'e--Birkhoff--Witt theorem for Clifford algebras~\cite{Chevalley1954}).
\end{definition}

\begin{remark}[Sign conventions]
For real signature $(p, q)$ ($n = p + q$), we write $\Cl_{p,q}(K)$ for the Clifford algebra of $Q(x_1, \dots, x_n) = \sum_{i=1}^p x_i^2 - \sum_{j=1}^q x_{p+j}^2$ under~\eqref{eq:cliff-rel}. Authors in Riemannian spin geometry~\cite{LawsonMichelsohn1989} often adopt the opposite sign $v^2 = -Q(v)\cdot 1$, which simply interchanges $(p, q) \leftrightarrow (q, p)$ (replacing $p - q \bmod 8$ by $q - p \bmod 8$). None of our positive or negative conclusions depend on this convention: $\Cl_{m,m}$ and $\Lambda(V)$ are sign-symmetric, odd dimensions remain odd, and in dimension $n = 2$ the division algebra $\mathbb{H}$ appears either at $\Cl_{0,2}$ (under $v^2 = +Q(v)$) or at $\Cl_{2,0}$ (under $v^2 = -Q(v)$).
\end{remark}

\begin{definition}[Minimal left ideals and algebraic spinors]
Let $A$ be a unital associative $K$-algebra. A $K$-subspace $S \subseteq A$ is a \emph{minimal (nonzero) left ideal} if:
\begin{enumerate}
  \item $A S \subseteq S$ ($S$ is a left ideal, hence a left $A$-module under left multiplication $\rho_S(a)(x) = ax$),
  \item $S \neq \{0\}$, and
  \item for every left ideal $I \subseteq S$, either $I = \{0\}$ or $I = S$ (equivalently, $S$ is a simple left $A$-module embedded as a left ideal of $A$).
\end{enumerate}
Following Cartan and Chevalley~\cite{Chevalley1954,Lounesto2001}, \emph{algebraic spinors} of $(V, Q)$ are elements of a minimal left ideal $S \subseteq \Cl(V, Q)$.
\end{definition}

\section{Positive Non-Vacuity Theorem: The Even-Dimensional Split Regime}

We first establish that in the classical Cartan--Chevalley setting---an even-dimensional nondegenerate split (neutral) quadratic space---every clause of Conjecture \texttt{00000007604} holds rigorously.

\begin{theorem}[Positive realization in even-dimensional split Clifford algebras]\label{thm:positive-split}
Let $K$ be a field with $\mathrm{char}\,K \neq 2$, and let $(V, Q)$ be a nondegenerate split quadratic space of even dimension $n = 2m$ ($m \ge 1$), so that $V = W \oplus W^*$ is the direct sum of $m$-dimensional totally isotropic subspaces with dual bases $(f_1, \dots, f_m)$ of $W$ and $(f_1^*, \dots, f_m^*)$ of $W^*$ satisfying
\[
  f_i^2 = (f_i^*)^2 = 0, \qquad f_i f_j + f_j f_i = 0, \qquad f_i^* f_j^* + f_j^* f_i^* = 0, \qquad f_i f_j^* + f_j^* f_i = \delta_{ij}\cdot 1.
\]
Set $e = (f_1^* f_1)\cdots(f_m^* f_m) \in \Cl(V, Q)$ and $S = \Cl(V, Q) e$. Then:
\begin{enumerate}
  \item \textbf{Minimal left ideal and exact half-dimension power of two:} $e^2 = e \neq 0$ is a primitive idempotent, $S = \Lambda(W) e$ is a minimal nonzero left ideal of $\Cl(V, Q)$ with ordered $K$-basis $\{f_{i_1}\cdots f_{i_k} e : 1 \le i_1 < \dots < i_k \le m\}$, and
  \[
    \dim_K S = 2^m = 2^{n/2}.
  \]
  \item \textbf{Faithful matrix isomorphism:} The left multiplication map $\rho_S : \Cl(V, Q) \to \End_K(S)$ is an isomorphism of $K$-algebras:
  \[
    \Cl(V, Q) \xrightarrow{\;\sim\;} \End_K(S) \cong M_{2^m}(K).
  \]
  Consequently, $\Cl(V, Q)$ is central simple and $S$ is (up to isomorphism) its unique simple left module.
  \item \textbf{Complete Cartan weight spectrum:} The elements $\sigma_j = f_j f_j^* - f_j^* f_j \in \Cl(V, Q)$ ($1 \le j \le m$) pairwise commute and satisfy $\sigma_j^2 = 1$. On $S$, each basis vector $f_I e$ ($I \subseteq \{1,\dots,m\}$) is a joint eigenvector with $\sigma_j (f_I e) = \varepsilon_{I,j} (f_I e)$, where $\varepsilon_{I,j} = +1$ if $j \in I$ and $\varepsilon_{I,j} = -1$ if $j \notin I$. Thus every sign tuple in $\{\pm 1\}^m$ occurs on $S$ with multiplicity~$1$.
  \item \textbf{Single-column matrix carrier:} Ordering the basis $\{b_1, \dots, b_{2^m}\}$ of $S$ with $b_1 = e = f_\emptyset e$, the induced matrix isomorphism $\Phi : \Cl(V, Q) \xrightarrow{\sim} M_{2^m}(K)$ satisfies $\Phi(e) = E_{1,1}$, and $\Phi(S) = M_{2^m}(K) E_{1,1}$ is the first-column subspace of $M_{2^m}(K)$, carrying the canonical left action of $M_{2^m}(K)$ on column vectors $K^{2^m}$.
\end{enumerate}
\end{theorem}

\begin{proof}
For each $j \in \{1,\dots,m\}$, set $e_j = f_j^* f_j$. Since $(f_j^*)^2 = 0$ and $f_j f_j^* = 1 - f_j^* f_j$, we have $e_j^2 = f_j^* (f_j f_j^*) f_j = f_j^*(1 - f_j^* f_j) f_j = f_j^* f_j = e_j$. For $i \neq j$, $f_i$ and $f_i^*$ anticommute with both $f_j$ and $f_j^*$, so $e_i$ commutes with $e_j$. Hence $e = e_1 \cdots e_m$ satisfies $e^2 = e$.

Using $f_j^* e_j = (f_j^*)^2 f_j = 0$ and commutativity of $e_j$ with $e_k$ ($k \neq j$), we obtain $f_j^* e = 0$ for all $j$. Every monomial in $\Cl(V, Q)$ can be reordered via $f_i f_j^* = \delta_{ij} - f_j^* f_i$ so that all $f^*$-factors stand to the right of all $f$-factors; acting on $e$, every term with at least one $f_j^*$ vanishes, leaving $S = \Cl(V, Q) e = \Span_K\{f_I e : I \subseteq \{1,\dots,m\}\}$, so $\dim_K S \le 2^m$. Moreover, for $I = \{i_1 < \dots < i_k\}$, multiplying $\sum_J c_J f_J e = 0$ on the left by $f_{i_k}^* \cdots f_{i_1}^*$ and projecting via $e$ isolates $c_I e = 0$; since in the exterior model on $\Lambda(W)$ (where $f_j$ acts by $f_j \wedge -$ and $f_j^*$ acts by interior contraction $\iota_{f_j^*}$) the operator corresponding to $e$ sends $1 \mapsto 1 \neq 0$, we have $e \neq 0$, linear independence of $\{f_I e\}$, and $\dim_K S = 2^m = 2^{n/2}$.

The interior/exterior action on $\Lambda(W)$ defines a $K$-algebra homomorphism $\Cl(V, Q) \to \End_K(\Lambda(W))$ intertwined with $\rho_S$ via $u \mapsto u \cdot e$ ($\Lambda(W) \xrightarrow{\sim} S$). For any subset $I, J \subseteq \{1,\dots,m\}$, the operator $\bigl(\prod_{i \in I} f_i\bigr) e \bigl(\prod_{j \in J}^{\text{rev}} f_j^*\bigr)$ sends $f_J e \mapsto f_I e$ and annihilates $f_K e$ for $K \neq J$, exhibiting all $2^m \times 2^m$ matrix units of $\End_K(S)$. Hence $\rho_S : \Cl(V, Q) \to \End_K(S) \cong M_{2^m}(K)$ is surjective; since $\dim_K \Cl(V, Q) = 2^{2m} = (2^m)^2 = \dim_K M_{2^m}(K)$, $\rho_S$ is an isomorphism. Minimality of $S = \Cl(V, Q) e$ and uniqueness of the simple module now follow immediately from the standard Wedderburn structure of $M_{2^m}(K)$, where $\Phi(e) = E_{1,1}$ identifies $\Phi(S) = M_{2^m}(K) E_{1,1}$ with the first column of $M_{2^m}(K)$.

Finally, $\sigma_j = f_j f_j^* - f_j^* f_j = 1 - 2 e_j$ commutes with $\sigma_k$ ($k \neq j$) and satisfies $\sigma_j^2 = 1 - 4e_j + 4e_j^2 = 1$. If $j \notin I$, then $e_j f_I e = f_I e_j e = f_I e$, so $\sigma_j(f_I e) = (1 - 2)(f_I e) = -f_I e$; if $j \in I$, writing $f_I = \pm f_j f_{I \setminus \{j\}}$ gives $e_j f_j = f_j^* f_j^2 = 0$, so $e_j(f_I e) = 0$ and $\sigma_j(f_I e) = +f_I e$. Thus the $2^m$ basis spinors $\{f_I e\}$ diagonalize $(\sigma_1,\dots,\sigma_m)$ with distinct eigenvalues covering all of $\{\pm 1\}^m$.
\end{proof}

\begin{example}[Concrete $n = 2$ witness $\Cl_{1,1}(\mathbb{Q}) \cong M_2(\mathbb{Q})$, formalized in Lean~4]\label{ex:cl11}
For $V = \mathbb{Q}^2$ with $Q(x_1, x_2) = x_1^2 - x_2^2$ ($e_1^2 = +1$, $e_2^2 = -1$, $e_{12} = e_1 e_2 = -e_2 e_1$), the map
\[
  \Phi_{1,1}(a_0 + a_1 e_1 + a_2 e_2 + a_{12} e_{12}) =
  \begin{pmatrix}
    a_0 + a_1 & -a_2 - a_{12} \\
    a_2 - a_{12} & a_0 - a_1
  \end{pmatrix}
\]
is an explicit $\mathbb{Q}$-algebra isomorphism $\Cl_{1,1}(\mathbb{Q}) \xrightarrow{\sim} M_2(\mathbb{Q})$. The primitive idempotent $e = \frac{1 + e_1}{2}$ satisfies $\Phi_{1,1}(e) = E_{1,1} = \bigl(\begin{smallmatrix}1&0\\0&0\end{smallmatrix}\bigr)$, and $S_{1,1} = \Cl_{1,1}(\mathbb{Q}) e$ has $\mathbb{Q}$-basis $s_+ = e$ and $s_- = e_2 e = \frac{e_2 - e_{12}}{2}$ with $\Phi_{1,1}(s_+) = \bigl(\begin{smallmatrix}1&0\\0&0\end{smallmatrix}\bigr)$, $\Phi_{1,1}(s_-) = \bigl(\begin{smallmatrix}0&0\\1&0\end{smallmatrix}\bigr)$, $\dim_{\mathbb{Q}} S_{1,1} = 2 = 2^{2/2}$, $e_1 s_\pm = \pm s_\pm$, and $\Phi_{1,1}(S_{1,1}) = \bigl\{\bigl(\begin{smallmatrix}u&0\\v&0\end{smallmatrix}\bigr) : u,v \in \mathbb{Q}\bigr\}$.
\end{example}

\begin{corollary}[Complex even dimensions]\label{cor:complex-even}
Over any algebraically closed field $K$ with $\mathrm{char}\,K \neq 2$ (in particular $K = \mathbb{C}$), every nondegenerate quadratic form $Q$ in even dimension $n = 2m$ is split ($\CCl_{2m} \cong M_{2^m}(\mathbb{C})$), so Theorem~\ref{thm:positive-split} holds for every nondegenerate even-dimensional complex Clifford algebra.
\end{corollary}

\section{Refutation Regime I: Odd Dimensions \texorpdfstring{($n = 2m + 1$)}{(n = 2m + 1)}}

We now show that Conjecture \texttt{00000007604} fails unconditionally in \emph{every} odd dimension $n = 2m + 1$, even over $\mathbb{C}$ or for split signatures over $\mathbb{Q}$ and $\mathbb{R}$.

\begin{proposition}[Arithmetic impossibility in odd dimensions]\label{prop:odd-arith}
For every integer $m \ge 0$, let $n = 2m + 1$. There exists no integer $d \in \mathbb{Z}$ and no rational number $q \in \mathbb{Q}$ satisfying $q^2 = 2^n$. In particular, no finite-dimensional vector space $S$ over any field $K$ can satisfy $\dim_K S = 2^{n/2}$ (nor $(\dim_K S)^2 = 2^n$).
\end{proposition}

\begin{proof}
Suppose $(a/b)^2 = 2^{2m+1}$ with $a, b \in \mathbb{Z}$, $b > 0$, $\gcd(a, b) = 1$. Then $a^2 = 2^{2m+1} b^2$; comparing $2$-adic valuations gives $2 v_2(a) = 2m + 1 + 2 v_2(b)$, reducing modulo~$2$ to $0 \equiv 1 \pmod{2}$, a contradiction.
\end{proof}

Beyond the arithmetic obstruction on $2^{n/2}$, odd-dimensional split/complex Clifford algebras fail the faithfulness, uniqueness, and spectral completeness clauses as well.

\begin{theorem}[Structural failure in odd-dimensional split and complex Clifford algebras]\label{thm:odd-split}
Let $K$ be a field with $\mathrm{char}\,K \neq 2$, and let $(V, Q)$ have odd dimension $n = 2m + 1$ with orthogonal basis $(e_1, \dots, e_{m+1}, u_1, \dots, u_m)$ satisfying $Q(e_i) = +1$ and $Q(u_j) = -1$ (giving $\Cl_{m+1,m}(K)$; over $K = \mathbb{C}$, every nondegenerate odd-dimensional space $\CCl_{2m+1}$ is isomorphic to this algebra). Define the signed volume element
\[
  \omega = e_{m+1} (e_1 u_1)\cdots(e_m u_m) \in \Cl_{m+1,m}(K).
\]
Then:
\begin{enumerate}
  \item $\omega$ lies in the center $Z(\Cl_{m+1,m}(K))$ and satisfies $\omega^2 = +1$.
  \item The central orthogonal idempotents $p_+ = \frac{1 + \omega}{2}$ and $p_- = \frac{1 - \omega}{2}$ decompose the algebra into two two-sided simple ideals:
  \[
    \Cl_{m+1,m}(K) = p_+ \Cl_{m+1,m}(K) \oplus p_- \Cl_{m+1,m}(K) \cong M_{2^m}(K) \oplus M_{2^m}(K).
  \]
  \item Every minimal nonzero left ideal $S \subset \Cl_{m+1,m}(K)$ lies either entirely in $p_+ \Cl_{m+1,m}(K)$ or entirely in $p_- \Cl_{m+1,m}(K)$, and satisfies:
  \begin{enumerate}
    \item $\dim_K S = 2^m = 2^{(n-1)/2} \neq 2^{n/2}$, with $(\dim_K S)^2 = 2^{2m} = 2^{n-1} \neq 2^n$;
    \item $S$ is \textbf{not faithful}: the complementary simple block $p_\mp \Cl_{m+1,m}(K)$ of dimension $2^{2m} = 2^{n-1} \ge 1$ annihilates $S$, so $\rho_S : \Cl_{m+1,m}(K) \to \End_K(S) \cong M_{2^m}(K)$ has nonzero kernel;
    \item \textbf{Incomplete spectrum:} Although $\omega \in \Cl_{m+1,m}(K)$ has eigenvalues $\{+1, -1\}$ in the regular representation, on any single minimal left ideal $S$ the operator $\omega|_S$ equals either $+\id_S$ or $-\id_S$, missing half of its spectrum (and $\Cl_{m+1,m}(K)$ has \emph{two} non-isomorphic simple modules $S_+$ and $S_-$, rather than ``the simple module'').
  \end{enumerate}
\end{enumerate}
\end{theorem}

\begin{proof}
Each pair $\sigma_j = e_j u_j$ ($1 \le j \le m$) satisfies $\sigma_j^2 = -e_j^2 u_j^2 = -(+1)(-1) = +1$ and commutes with $e_{m+1}$ as well as with $\sigma_k$ for $k \neq j$. Furthermore, $e_j \sigma_j = e_j^2 u_j = u_j = -\sigma_j e_j$ and $u_j \sigma_j = -e_j u_j^2 = e_j = -\sigma_j u_j$, so both $e_j$ and $u_j$ anticommute with $\sigma_j$ and also anticommute with $e_{m+1}$ (while commuting with $\sigma_k$ for $k \neq j$). Consequently, conjugating $\omega = e_{m+1}\sigma_1\cdots\sigma_m$ by $e_j$ or $u_j$ produces two sign reversals that cancel, and conjugating by $e_{m+1}$ produces none; hence $\omega \in Z(\Cl_{m+1,m}(K))$. Moreover,
\[
  \omega^2 = e_{m+1}^2 \sigma_1^2 \cdots \sigma_m^2 = (+1)(+1)^m = +1.
\]
Therefore $p_\pm = \frac{1 \pm \omega}{2}$ are nonzero orthogonal central idempotents ($p_+ + p_- = 1$, $p_+ p_- = 0$). Restricting the even subalgebra $\Cl_{m,m}(K) = \langle e_1,\dots,e_m, u_1,\dots,u_m \rangle \cong M_{2^m}(K)$ via $x \mapsto p_\pm x$ yields $K$-algebra isomorphisms $\Cl_{m,m}(K) \xrightarrow{\sim} p_\pm \Cl_{m+1,m}(K) \cong M_{2^m}(K)$ (injective by simplicity of $M_{2^m}(K)$ and surjective by dimension counting $2^{2m} + 2^{2m} = 2^{2m+1}$).

Now let $S \subset \Cl_{m+1,m}(K)$ be any minimal nonzero left ideal. Since $S = p_+ S \oplus p_- S$ is a direct sum of left ideals inside $S$, minimality forces either $p_- S = 0$ (so $S \subseteq p_+ \Cl_{m+1,m}(K) \cong M_{2^m}(K)$) or $p_+ S = 0$ (so $S \subseteq p_- \Cl_{m+1,m}(K) \cong M_{2^m}(K)$). In either case, $S$ is a minimal left ideal of $M_{2^m}(K)$, hence $\dim_K S = 2^m = 2^{(n-1)/2}$, $\ker \rho_S = p_\mp \Cl_{m+1,m}(K)$ has dimension $2^{2m} > 0$, and $\omega$ acts on $S$ purely by the scalar $+1$ (if $S \subset p_+ \Cl$) or $-1$ (if $S \subset p_- \Cl$).
\end{proof}

\begin{corollary}[Minimal counterexample at $n = 1$: $\Cl_{1,0}(K) \cong K \times K$]\label{cor:cl10}
At $n = 1$ ($m = 0$) with $Q(x) = x^2$, $\Cl_{1,0}(K) = K \oplus K e_1$ ($e_1^2 = 1$) has exactly two minimal nonzero left ideals,
\[
  S_+ = K\Bigl(\frac{1 + e_1}{2}\Bigr), \qquad S_- = K\Bigl(\frac{1 - e_1}{2}\Bigr),
\]
each of dimension $\dim_K S_\pm = 1$ (so $(\dim_K S_\pm)^2 = 1 \neq 2 = 2^1$), with nonzero annihilator $\ker \rho_{S_\pm} = S_\mp \neq \{0\}$ and incomplete generator spectrum $\mathrm{Spec}(e_1|_{S_\pm}) = \{\pm 1\} \subsetneq \{+1, -1\}$.
\end{corollary}

\section{Refutation Regime II: Quaternionic and Bott Obstructions}

One might conjecture that restricting to \emph{even} dimension $n = 2m$ and \emph{nondegenerate} quadratic forms $Q$ rescues Conjecture \texttt{00000007604}. Over non-algebraically-closed fields---including $\mathbb{Q}$ and the physically foundational real field $\mathbb{R}$---this fails because a central simple algebra over $K$ has the Wedderburn form $M_k(\mathbb{D})$ for a central division $K$-algebra $\mathbb{D}$, which need not be $K$ itself.

\begin{proposition}[General Brauer/Wedderburn dimension inflation]\label{prop:wedderburn}
Suppose $\dim_K V = n = 2m$ is even, $Q$ is nondegenerate, and $\Cl(V, Q) \cong M_k(\mathbb{D})$ where $\mathbb{D}$ is a central division $K$-algebra of Schur index $d \ge 1$ ($\dim_K \mathbb{D} = d^2$). Since $2^{2m} = \dim_K \Cl(V, Q) = k^2 d^2$, we have $k d = 2^m = 2^{n/2}$, and every minimal nonzero left ideal $S \subset \Cl(V, Q)$ is isomorphic to the column module $\mathbb{D}^k$, whose $K$-dimension is
\begin{equation}\label{eq:schur-dim}
  \dim_K S = k \cdot \dim_K \mathbb{D} = k d^2 = d \cdot (kd) = d \cdot 2^{n/2}.
\end{equation}
Thus $\dim_K S = 2^{n/2}$ holds if and only if $d = 1$ ($\mathbb{D} = K$, i.e.\ $\Cl(V, Q)$ is split). Whenever $d > 1$, $\dim_K S = d \cdot 2^{n/2} > 2^{n/2}$, and the image $\rho_S(\Cl(V, Q)) \cong M_k(\mathbb{D})$ has $K$-dimension $2^n$, which is a proper subalgebra of $\End_K(S) \cong M_{d \cdot 2^{n/2}}(K)$ (of dimension $d^2 2^n > 2^n$, with nontrivial commutant $\End_{\Cl}(S) \cong \mathbb{D}^{\mathrm{op}}$).
\end{proposition}

\begin{theorem}[Even-dimensional counterexample $\Cl_{0,2}(K) \cong \mathbb{H}_K$ at $n = 2$]\label{thm:cl02}
Let $K$ be any formally real field (in particular $K = \mathbb{Q}$ or $K = \mathbb{R}$), and let $V = K^2$ ($n = 2$) carry the nondegenerate negative-definite quadratic form $Q(x_1, x_2) = -x_1^2 - x_2^2$. Let $\Cl_{0,2}(K)$ have orthogonal generators $e_1, e_2$ with $e_1^2 = -1$, $e_2^2 = -1$, and $e_{12} = e_1 e_2 = -e_2 e_1$ ($e_{12}^2 = -1$). Then:
\begin{enumerate}
  \item \textbf{Quaternion division ring:} For every $q = a_0 + a_1 e_1 + a_2 e_2 + a_{12} e_{12} \in \Cl_{0,2}(K)$, set $\bar{q} = a_0 - a_1 e_1 - a_2 e_2 - a_{12} e_{12}$ and $N(q) = a_0^2 + a_1^2 + a_2^2 + a_{12}^2 \in K$. Then $q \bar{q} = \bar{q} q = N(q) \cdot 1$, and because $K$ is formally real, $N(q) = 0 \iff q = 0$. Hence every $q \neq 0$ has two-sided inverse $q^{-1} = N(q)^{-1}\bar{q}$, so $\Cl_{0,2}(K) \cong \mathbb{H}_K$ is a division ring.
  \item \textbf{Absence of proper idempotents and unique minimal left ideal of dimension $4 \neq 2$:} The only idempotents in $\Cl_{0,2}(K)$ are $0$ and $1$, and the \textbf{unique} minimal nonzero left ideal of $\Cl_{0,2}(K)$ is $S = \Cl_{0,2}(K)$ itself, with
  \[
    \dim_K S = 4 = 2 \cdot 2^{2/2} \neq 2 = 2^{2/2}.
  \]
  \item \textbf{Failure of full matrix realization and real/rational spectrum:} Left multiplication embeds $\Cl_{0,2}(K)$ as a $4$-dimensional subalgebra of $\End_K(S) \cong M_4(K)$ ($\dim_K M_4(K) = 16$), not $M_2(K)$, and no nonzero vector $v \in V \setminus \{0\}$ admits any eigenvalue $\lambda \in K$ on $S$.
\end{enumerate}
\end{theorem}

\begin{proof}
Direct expansion in the basis $(1, e_1, e_2, e_{12})$ gives $q \bar{q} = \bar{q} q = (a_0^2 + a_1^2 + a_2^2 + a_{12}^2)\cdot 1$ because all cross terms cancel via $e_1 e_2 = -e_2 e_1$, $e_1 e_{12} = -e_2 = -e_{12} e_1$, and $e_2 e_{12} = e_1 = -e_{12} e_2$. Over a formally real field $K$, a sum of squares vanishes if and only if each summand vanishes, so $q \neq 0 \implies N(q) \neq 0$ and $q (N(q)^{-1}\bar{q}) = (N(q)^{-1}\bar{q}) q = 1$.

If $e^2 = e$ and $e \neq 0$, multiplying on the left by $e^{-1}$ yields $e = 1$. If $I \subseteq \Cl_{0,2}(K)$ is any nonzero left ideal containing $q \neq 0$, then for every $x \in \Cl_{0,2}(K)$, $x = (x q^{-1}) q \in I$, so $I = \Cl_{0,2}(K)$. Thus $S = \Cl_{0,2}(K)$ is the unique minimal nonzero left ideal, of dimension $\dim_K S = 4 \neq 2$. Finally, if $v \in V \setminus \{0\}$ satisfied $v s = \lambda s$ for some $s \in S \setminus \{0\}$ and $\lambda \in K$, multiplying on the left by $v$ and on the right by $s^{-1}$ would give $Q(v)\cdot 1 = v^2 = \lambda^2 \cdot 1$, i.e.\ $\lambda^2 + v_1^2 + v_2^2 = 0$ with $(v_1, v_2) \neq (0,0)$, impossible in $K$.
\end{proof}

Combining Theorem~\ref{thm:positive-split}, Theorem~\ref{thm:odd-split}, and Proposition~\ref{prop:wedderburn} with the Atiyah--Bott--Shapiro classification~\cite{AtiyahBottShapiro1964} yields the complete modulo-$8$ determination of real minimal spinor dimensions.

\begin{theorem}[Complete real Bott periodicity classification of minimal left ideal dimensions]\label{thm:bott}
Let $n = p + q \ge 1$ and write $n = 2m$ (if even) or $n = 2m + 1$ (if odd). Under $v^2 = +Q(v)\cdot 1$, the structure of $\Cl_{p,q}(\mathbb{R})$ and the real dimension $\dim_{\mathbb{R}} S$ of every minimal nonzero left ideal $S \subset \Cl_{p,q}(\mathbb{R})$ depend solely on $n$ and $s = (p - q) \bmod 8$ according to Table~\ref{tab:bott}. In particular, for $n \ge 1$,
\[
  \dim_{\mathbb{R}} S = 2^{n/2} \quad \Longleftrightarrow \quad n \text{ is even and } p - q \equiv 0 \text{ or } 2 \pmod{8}.
\]
\end{theorem}

\begin{table}[ht]
\centering
\footnotesize
\setlength{\tabcolsep}{4.5pt}
\renewcommand{\arraystretch}{1.22}
\begin{tabular}{@{}cccccl@{}}
\hline
$(p-q)\bmod 8$ & Parity of $n$ & $\Cl_{p,q}(\mathbb{R})$ & Simple Ideal $S$ & $\dim_{\mathbb{R}} S$ & Comparison with $2^{n/2}$ \\
\hline
$0$ & $n = 2m$ & $M_{2^m}(\mathbb{R})$ & $\mathbb{R}^{2^m}$ & $2^m$ & $= 2^{n/2}$ (\textbf{holds}) \\
$1$ & $n = 2m{+}1$ & $M_{2^m}(\mathbb{R}) \oplus M_{2^m}(\mathbb{R})$ & $\mathbb{R}^{2^m}$ & $2^m$ & $= 2^{(n-1)/2} \neq 2^{n/2}$ (fails) \\
$2$ & $n = 2m$ & $M_{2^m}(\mathbb{R})$ & $\mathbb{R}^{2^m}$ & $2^m$ & $= 2^{n/2}$ (\textbf{holds}) \\
$3$ & $n = 2m{+}1$ & $M_{2^m}(\mathbb{C})$ & $\mathbb{C}^{2^m}$ & $2^{m+1}$ & $= 2^{(n+1)/2} \neq 2^{n/2}$ (fails) \\
$4$ & $n = 2m$ & $M_{2^{m-1}}(\mathbb{H})$ & $\mathbb{H}^{2^{m-1}}$ & $2^{m+1}$ & $= 2 \cdot 2^{n/2} \neq 2^{n/2}$ (fails) \\
$5$ & $n = 2m{+}1$ & $M_{2^{m-1}}(\mathbb{H}) \oplus M_{2^{m-1}}(\mathbb{H})$ & $\mathbb{H}^{2^{m-1}}$ & $2^{m+1}$ & $= 2^{(n+1)/2} \neq 2^{n/2}$ (fails) \\
$6$ & $n = 2m$ & $M_{2^{m-1}}(\mathbb{H})$ & $\mathbb{H}^{2^{m-1}}$ & $2^{m+1}$ & $= 2 \cdot 2^{n/2} \neq 2^{n/2}$ (fails) \\
$7$ & $n = 2m{+}1$ & $M_{2^m}(\mathbb{C})$ & $\mathbb{C}^{2^m}$ & $2^{m+1}$ & $= 2^{(n+1)/2} \neq 2^{n/2}$ (fails) \\
\hline
\end{tabular}
\caption{Modulo-$8$ periodicity of $\Cl_{p,q}(\mathbb{R})$ ($n = p + q \ge 1$) and real minimal left ideal dimension $\dim_{\mathbb{R}} S$ under $v^2 = +Q(v)\cdot 1$ (for $s \in \{4,5,6\}$, $n \ge 2$ so $m \ge 1$).}
\label{tab:bott}
\end{table}

\begin{proof}
The algebra isomorphisms in column~3 are the classical periodicity relations $\Cl_{1,1}(\mathbb{R}) \cong \Cl_{2,0}(\mathbb{R}) \cong M_2(\mathbb{R})$, $\Cl_{0,2}(\mathbb{R}) \cong \mathbb{H}$, $\Cl_{1,0}(\mathbb{R}) \cong \mathbb{R} \oplus \mathbb{R}$, $\Cl_{0,1}(\mathbb{R}) \cong \mathbb{C}$, $\Cl_{p+1,q+1}(\mathbb{R}) \cong M_2(\Cl_{p,q}(\mathbb{R}))$, $\Cl_{p+2,q}(\mathbb{R}) \cong \Cl_{q,p}(\mathbb{R}) \otimes_{\mathbb{R}} M_2(\mathbb{R})$, and $\Cl_{p,q+2}(\mathbb{R}) \cong \Cl_{q,p}(\mathbb{R}) \otimes_{\mathbb{R}} \mathbb{H}$~\cite{AtiyahBottShapiro1964,LawsonMichelsohn1989}. For $M_k(\mathbb{D})$ (or $M_k(\mathbb{D}) \oplus M_k(\mathbb{D})$), every minimal nonzero left ideal is isomorphic as a real vector space to $\mathbb{D}^k$, of real dimension $k \cdot \dim_{\mathbb{R}}\mathbb{D}$ where $\dim_{\mathbb{R}} \mathbb{R} = 1$, $\dim_{\mathbb{R}} \mathbb{C} = 2$, and $\dim_{\mathbb{R}} \mathbb{H} = 4$. Comparing column~5 with $2^{n/2}$ establishes the stated equivalence.
\end{proof}

\section{Refutation Regime III: Degenerate Quadratic Forms and Exterior Algebras}

Finally, even over an algebraically closed field such as $K = \mathbb{C}$ and in even dimension $n = 2m \ge 2$, Conjecture \texttt{00000007604} fails catastrophically whenever the quadratic form $Q$ is degenerate.

\begin{theorem}[Collapse of minimal left ideals in the exterior Clifford algebra $\Lambda(V)$]\label{thm:exterior}
Let $K$ be any field, let $V$ be a $K$-vector space of dimension $n \ge 1$ with basis $(e_1, \dots, e_n)$, and let $Q = 0$, so that $\Cl(V, 0) = \Lambda(V) = \bigoplus_{k=0}^n \Lambda^k(V)$ is the exterior algebra. Then:
\begin{enumerate}
  \item \textbf{Local non-semisimple structure:} The augmentation ideal $\mathfrak{m} = \bigoplus_{k=1}^n \Lambda^k(V)$ of dimension $2^n - 1$ is nilpotent ($\mathfrak{m}^{n+1} = 0$) and equals the Jacobson radical $J(\Lambda(V))$. Every $a \in \Lambda(V) \setminus \mathfrak{m}$ is a unit, so $\Lambda(V)$ is a local Artinian algebra whose only idempotents are $0$ and $1$.
  \item \textbf{Unique $1$-dimensional minimal left ideal:} Every nonzero left ideal $I \subseteq \Lambda(V)$ contains the top-degree line $\Lambda^n(V) = K(e_1 \wedge \dots \wedge e_n)$. Consequently, $S_{\mathrm{top}} = \Lambda^n(V)$ is the \textbf{unique} minimal nonzero left ideal of $\Lambda(V)$, and
  \[
    \dim_K S_{\mathrm{top}} = 1 \neq 2^{n/2} \qquad \text{for all } n \ge 1
  \]
  (in particular, in even dimension $n = 2m \ge 2$, $\dim_K S_{\mathrm{top}} = 1 < 2^m = 2^{n/2}$).
  \item \textbf{Nilpotence and trivial Clifford action:} $S_{\mathrm{top}}^2 = 0$ (so $S_{\mathrm{top}}$ contains no nonzero idempotent and cannot be written as $\Lambda(V) e$ for any idempotent $e$), and $v \wedge s = 0$ for all $v \in V$ and $s \in S_{\mathrm{top}}$. Thus $\rho_{S_{\mathrm{top}}} : \Lambda(V) \to \End_K(S_{\mathrm{top}}) \cong K$ is merely the scalar augmentation map with kernel $\mathfrak{m}$ of dimension $2^n - 1$, annihilating the entire generating space $V$.
\end{enumerate}
\end{theorem}

\begin{proof}
Every $x \in \mathfrak{m}$ is a sum of homogeneous elements of positive degree; since $\Lambda^k(V) = 0$ for $k > n$, any product of $n + 1$ elements of $\mathfrak{m}$ vanishes, so $x^{n+1} = 0$. If $a = c_0 \cdot 1 + x$ with $c_0 \in K^\times$ and $x \in \mathfrak{m}$, then $a^{-1} = c_0^{-1}\sum_{j=0}^n (-c_0^{-1}x)^j$ is a two-sided inverse. Hence $\mathfrak{m}$ is the unique maximal left (and two-sided) ideal, i.e.\ $J(\Lambda(V)) = \mathfrak{m}$, and any idempotent $e \in \Lambda(V)$ satisfies $e(1 - e) = 0$; since $e \in \mathfrak{m}$ or $1 - e \in \mathfrak{m}$ cannot both hold, one of $e, 1 - e$ is a unit, forcing $e \in \{0, 1\}$.

Next let $I \subseteq \Lambda(V)$ be any nonzero left ideal and pick $0 \neq y \in I$. Write $y = \sum_{J \subseteq \{1,\dots,n\}} c_J e_J$ in the ordered wedge basis $e_J = \bigwedge_{j \in J} e_j$, and choose a subset $J_0$ of \emph{minimal cardinality} $|J_0| = k$ such that $c_{J_0} \neq 0$. Set $J_0^c = \{1,\dots,n\} \setminus J_0$ of cardinality $n - k$. For every $J \subseteq \{1,\dots,n\}$ with $c_J \neq 0$:
\begin{itemize}
  \item if $|J| > k$, then $|J_0^c| + |J| > (n - k) + k = n$, so $e_{J_0^c} \wedge e_J = 0$;
  \item if $|J| = k$ and $J \neq J_0$, then $J \cap J_0^c \neq \emptyset$, so $e_{J_0^c} \wedge e_J = 0$;
  \item for $J = J_0$, $e_{J_0^c} \wedge e_{J_0} = \varepsilon \, e_{\{1,\dots,n\}}$ with $\varepsilon \in \{\pm 1\}$.
\end{itemize}
Therefore $(\varepsilon c_{J_0}^{-1} e_{J_0^c}) \wedge y = e_{\{1,\dots,n\}} \in I$, so $\Lambda^n(V) \subseteq I$. Since $\dim_K \Lambda^n(V) = 1$, $S_{\mathrm{top}} = \Lambda^n(V)$ is the unique minimal nonzero left ideal. Because $n \ge 1$, $2n > n$ implies $S_{\mathrm{top}}^2 \subseteq \Lambda^{2n}(V) = 0$, and $V \wedge \Lambda^n(V) \subseteq \Lambda^{n+1}(V) = 0$.
\end{proof}

\begin{remark}[General degenerate quadratic forms]
More generally, if $\rad(B_Q) = V^\perp = \{z \in V : B_Q(z, V) = 0 \text{ and } Q(z) = 0\}$ has dimension $r \ge 1$, then $V = V_0 \perp R$ with $R = V^\perp$ ($\dim_K R = r$) and $V_0$ nondegenerate of dimension $n - r$. Any nonzero $z \in R$ anticommutes with $V$ and satisfies $z^2 = 0$, so the two-sided ideal $(R) = \Cl(V, Q) R \Cl(V, Q)$ is nilpotent ($(R)^{r+1} = 0$) and equals $J(\Cl(V, Q))$, with $\Cl(V, Q)/J(\Cl(V, Q)) \cong \Cl(V_0, Q|_{V_0})$. Every minimal nonzero left ideal $S \subseteq \Cl(V, Q)$ is annihilated on the left by $J(\Cl(V, Q))$, hence lies inside the left socle $\Lambda^r(R) \cdot \Cl(V, Q)$ where $S^2 = 0$ and $R$ acts as zero; as a module over $\Cl(V_0, Q|_{V_0})$, $S$ is a minimal left ideal of a Clifford algebra in strictly smaller dimension $n - r < n$ (so over $\mathbb{C}$, $\dim_{\mathbb{C}} S = 2^{\lfloor (n-r)/2 \rfloor} < 2^{n/2}$).
\end{remark}

\section{Formal Verification in Core Lean 4.19.0 and Computational Suite}

\subsection{Core Lean 4.19.0 Architecture}
The core algebraic structures, positive split witness, and all three refutation regimes are formalized in \textbf{Core Lean~4.19.0} (\texttt{leanprover/lean4:v4.19.0}, \texttt{-DwarningAsError=true}) inside \texttt{lean/} with zero external dependencies, zero \texttt{sorry}/\texttt{admit}, zero \texttt{native\_decide}, and zero custom axioms:
\begin{itemize}
  \item \texttt{PolyNorm.lean}: Self-contained verified sparse multivariate polynomial normalizer over $\mathbb{Z}$ using pure structural pattern matching so polynomial identities reduce in the Lean kernel via \texttt{decide}.
  \item \texttt{RatField.lean}: Quotient construction of the exact rational field $\mathbb{Q}$ (\texttt{RatQ}), complete field axioms, four-square anisotropy (\texttt{four\_squares\_eq\_zero}), irrationality of $\sqrt{2^{2m+1}}$ for all $m \in \mathbb{N}$ (\texttt{odd\_pow\_two\_not\_nat\_sq} and \texttt{odd\_pow\_two\_not\_rat\_sq}), and direct lifting of \texttt{PolyNorm} to \texttt{RatQ}.
  \item \texttt{AlgebraicFramework.lean}: Unital associative $\mathbb{Q}$-algebras (\texttt{QAlgebra}), left and two-sided ideals, minimal nonzero left ideals (\texttt{IsMinimalLeftIdeal}), faithfulness, nilpotence, exact $\mathbb{Q}$-linear dimension \texttt{HasQDim} via explicit linear bijection with $\mathbb{Q}^d$, and unconditional dimension-invariance theorems (\texttt{dim\_one\_not\_dim\_two} and \texttt{dim\_four\_not\_dim\_two}).
  \item \texttt{PositiveSplitCl11.lean}: Parametric Clifford families \texttt{Cl1} ($n = 1$) and \texttt{Cl2} ($n = 2$) proving $v^2 = Q(v)\cdot 1$, the $2\times 2$ matrix algebra \texttt{M2Q}, the explicit isomorphism $\Phi_{1,1} : \Cl_{1,1}(\mathbb{Q}) \xrightarrow{\sim} M_2(\mathbb{Q})$ (\texttt{Cl11.isoM2}), and full positive verification of $S_{1,1} = \Cl_{1,1}(\mathbb{Q}) e$: minimality, dimension $2 = 2^{2/2}$, faithfulness, complete $\{+1, -1\}$ Cartan weight spectrum, and single-column matrix carrier in $M_2(\mathbb{Q})$.
  \item \texttt{Counterexamples.lean}: Complete classifications of \emph{all} minimal nonzero left ideals, proving the refutations in odd dimension $\Cl_{1,0}(\mathbb{Q})$ (\texttt{Cl10.refutation\_cl10}), even quaternionic signature $\Cl_{0,2}(\mathbb{Q}) \cong \mathbb{H}_{\mathbb{Q}}$ (\texttt{Cl02.refutation\_cl02}), and even degenerate signature $\Lambda(\mathbb{Q}^2)$ (\texttt{Ext2.refutation\_ext2}).
  \item \texttt{SpinorMinimalIdeals.lean}: Root module assembling the four-part master resolution theorem (\texttt{conjecture\_00000007604\_complete\_resolution}).
\end{itemize}

\subsection{Standalone Exact Computational Verification}
The standalone Python~3 script \texttt{verify.py} (using only \texttt{fractions.Fraction} from the standard library) independently constructs $\Cl(V, Q)$ over $\mathbb{Q}$ for diagonal forms of signature $(p, q, r)$ up to $n = 4$, verifies basis associativity on $24$ Clifford algebras, computes exact centers and Dickson trace Jacobson radicals, verifies the positive split realizations $\Cl_{1,1}$ ($\dim S = 2$) and $\Cl_{2,2}$ ($\dim S = 4$ with all $4$ joint Cartan sign weights), confirms the refutations in $\Cl_{1,0}$, $\Cl_{0,1}$, $\Cl_{2,1}$, $\Cl_{0,2} = \mathbb{H}$, $\Cl_{4,0}/\Cl_{0,4} \cong M_2(\mathbb{H})$, $\Lambda(V_1),\dots,\Lambda(V_4)$, and $\Cl(\mathrm{diag}(1,0))$, and checks all $44$ signatures $(p, q)$ with $1 \le p + q \le 8$ against the modulo-$8$ Bott periodicity classification of Theorem~\ref{thm:bott}.

\begin{thebibliography}{9}

\bibitem{AtiyahBottShapiro1964}
M.~F. Atiyah, R.~Bott, and A.~Shapiro,
\emph{Clifford modules},
Topology \textbf{3} (1964), Suppl.~1, 3--38.

\bibitem{Cartan1966}
\'E.~Cartan,
\emph{The Theory of Spinors},
Hermann, Paris, 1966 (English translation of \emph{Le\c{c}ons sur la th\'eorie des spineurs}, 1938).

\bibitem{Chevalley1954}
C.~Chevalley,
\emph{The Algebraic Theory of Spinors},
Columbia University Press, New York, 1954.

\bibitem{LawsonMichelsohn1989}
H.~B. Lawson, Jr. and M.-L. Michelsohn,
\emph{Spin Geometry},
Princeton Mathematical Series, vol.~38, Princeton University Press, Princeton, NJ, 1989.

\bibitem{Lounesto2001}
P.~Lounesto,
\emph{Clifford Algebras and Spinors}, 2nd ed.,
London Mathematical Society Lecture Note Series, vol.~286, Cambridge University Press, Cambridge, 2001.

\end{thebibliography}

\end{document}
